% ============================================================
%  CHAPTER 1 — INTRODUCTION
% ============================================================
\chapter{Introduction}
\label{chap:intro}

Governments worldwide depend on Labour Force Surveys (LFS) as their
foundational source of workforce intelligence, supplying the indicators that
drive minimum-wage legislation, unemployment insurance schemes, vocational
training investments, and bilateral labour migration agreements.
The underlying microdata flow into cross-national repositories operated by
the ILO and the OECD, anchoring the labour-market comparisons that
shape policy in economies representing billions of workers~\cite{ilo2012isco_v1,zhang2025synthetic}.

The multilingual GCC context provides a particularly compelling motivation for this thesis.
With a resident population comprising over 200 nationalities, Arabic, English,
Urdu, Hindi, and Tagalog are all in daily commercial and domestic use.
Traditional survey instruments force respondents into a single language per
session, suppressing the natural code-switching behaviour that characterises
multilingual communities and introducing systematic measurement error among
linguistic minorities~\cite{hamed2025codeswitching}.
Simultaneously, the LFS requires classification to three simultaneous
international standards --- ISCO-08 for occupation, ISIC Rev.~4 for industry,
and ISCED~2011 for education --- a triple-classification burden no prior
automated survey system has integrated into a single pipeline.

% ------------------------------------------------------------
\section{Motivation and Background}
\label{sec:motivation}
% ------------------------------------------------------------

Contemporary LFS administration faces four interrelated operational challenges.

\subsection{Data Completeness and Real-Time Consistency Enforcement}

Across major national programmes, between one in seven and one in ten
survey items are left unanswered~\cite{zhang2025synthetic}.
Beyond missing values, the ILO's 19th International Conference of Labour
Statisticians (ICLS-19) defines ten cross-answer consistency rules (R01--R10)
that must be enforced to produce valid labour-force status determinations.
These rules include: the wage-gate rule (E1) that skips monthly wage questions
for employers and the self-employed, and the re-routing rule (F3) that
automatically reclassifies a respondent to ``not in labour force'' when they
report neither active job search (F1~=~no) nor availability for work
(F3~=~no).
Paper-based and conventional CAPI instruments cannot enforce these rules
in real time, requiring costly post-collection editing~\cite{zhang2025synthetic}.

\subsection{Triple-Standard Classification Complexity}

The ISCO-08 taxonomy provides 436 unit groups at the four-digit level (the
implemented system covers 441, reflecting the latest ILO supplementary
guidance) organised into 130 minor groups, 43 sub-major groups, and 10 major
groups~\cite{ilo2012isco_v1,ilo2012isco_v2}.
When trained specialists independently assign codes to the same job description,
they reach the same four-digit ISCO-08 unit group in only three out of four attempts on average~\cite{hamed2025isco}.

However, occupation coding in isolation is insufficient: correct LFS analysis
requires simultaneous classification of the employing industry to ISIC Rev.~4
(Section $\rightarrow$ Division $\rightarrow$ Group $\rightarrow$ 4-digit Class,
e.g.\ J $\rightarrow$ 62 $\rightarrow$ 620 $\rightarrow$ 6201 ``Computer
programming activities'') and the respondent's education to both ISCED~2011
attainment level (0--8) and ISCED-F~2013 field of specialisation (Broad
$\rightarrow$ Narrow $\rightarrow$ 4-digit Detailed, e.g.\ 06 $\rightarrow$
061 $\rightarrow$ 0613 ``Software and applications development'').
No prior conversational system has produced all three classifications within a
single survey session~\cite{hamed2025isco}.

\subsection{Multilingual and Arabic Dialect Complexity}

In multilingual GCC contexts, respondents naturally interleave Arabic (including
forms), English, Urdu, Hindi, and Tagalog within a single conversational turn.
Dialect normalisation requires approximately 30 systematic
dialect-to-MSA token replacements before semantic embedding can produce
reliable recall --- for example, the dialect negation \textit{shno} to MSA
\textit{ma}, or \textit{brouh} (first-person future) to \textit{adhhab}
--- a requirement absent from all prior Arabic LFS NLP
work~\cite{hamed2025codeswitching}.

\subsection{Resource Intensity and Operational Cost}

The door-to-door collection model routinely costs more than USD~100 per completed record~\cite{zhang2025synthetic}.
The Person Register pre-fill pattern --- using data from previous survey rounds
stored in the \texttt{person\_register} database table to verify rather than
re-collect confirmed fields --- reduces the active question count by 40--50\%,
directly reducing session duration and interviewer cost.

\subsection{Summary of Challenges}

\begin{table}[htbp]
\centering
\caption{Operational challenges in traditional LFS administration and their
documented impacts.}
\label{tab:ch1-challenges}
\small\renewcommand{\arraystretch}{1.35}
\resizebox{\textwidth}{!}{%
\begin{tabular}{>{\raggedright}p{3.0cm}
                >{\raggedright}p{5.6cm}
                >{\raggedright\arraybackslash}p{3.6cm}}
\toprule
\textbf{Challenge} &
\textbf{Documented\_Impact} &
\textbf{Evidence} \\
\midrule
High operational cost & Exceeds USD~100/interview; constrains frequency
  and geographic coverage & \cite{zhang2025synthetic} \\
Data quality deficits & 10--15\% item nonresponse; ILO ICLS-19 rules
  violated post-hoc; costly editing & \cite{zhang2025synthetic} \\
Multi-standard coding & 75--85\% inter-coder agreement; no prior system
  codes ISCO + ISIC + ISCED simultaneously &
  \cite{hamed2025isco} \\
Multilingual complexity & Code-switching produces measurement bias;
  regional Arabic dialect corrupts semantic embeddings &
  \cite{hamed2025codeswitching} \\
\bottomrule
\end{tabular}
}
\end{table}

% ------------------------------------------------------------
\section{The AI Opportunity}
\label{sec:opportunity}
% ------------------------------------------------------------

Large language models, multilingual embedding models, RAG frameworks, and
agentic orchestration together create a genuine opportunity to address all four
challenges simultaneously.
Key enablers in this thesis are:

\begin{itemize}
  \item \texttt{intfloat/\-multilingual-e5-large} (1024-dim sentence-transformers):
  provides robust semantic search across English, Arabic, Urdu, Hindi, and
  Tagalog without per-language fine-tuning~\cite{stump2025embedding}.

  \item \textbf{CrewAI hierarchical orchestration}: enables a 13-agent pipeline
  where specialist agents --- language detection, occupation coding, industry
  classification, education classification, validation, emotional intelligence,
  HITL management, audit logging, and report generation --- each operate with
  defined inputs, outputs, and escalation rules~\cite{crewai2024}.

  \item \textbf{Ollama + Llama~3.2} for general tasks (NER, conversation,
  summaries) with \textbf{Claude~3.5 Sonnet} fallback for critical tasks
  (ISCO re-ranking, validation, bilingual reports)~\cite{anthropic2024claude}.

  \item \textbf{Deterministic ILO/UNESCO correspondence tables} for the
  Semantic Relation Engine: the ILO Geneva 2012 ISCO-ISIC table and UNESCO
  ISCED~2011 Operational Manual Table~7 ensure cross-standard coherence
  checking is fully auditable without LLM involvement~\cite{ilo2012isco_v1,unesco2012isced}.
\end{itemize}

Deploying AI in official statistics imposes requirements absent from commercial
chatbot guidelines: full traceability of every classification decision;
strict version-pinning to specific standard revisions; calibrated confidence
with deterministic HITL escalation; tamper-evident GDPR-compliant audit
logging; and bilingual output in English and Arabic.

% ------------------------------------------------------------
\section{Research Objectives}
\label{sec:objectives}
% ------------------------------------------------------------

The primary aim is to design, implement, and empirically validate a multilingual
conversational agentic system for the complete Labour Force Survey that
simultaneously classifies respondents to ISCO-08, ISIC Rev.~4, and ISCED~2011,
enforces ILO ICLS-19 consistency rules in real time, and cross-validates the
three classifications through a novel Semantic Relation Engine.

\begin{table}[htbp]
\centering
\caption{Research objectives and target metrics.}
\label{tab:ch1-objectives}
\small\renewcommand{\arraystretch}{1.35}
\resizebox{\textwidth}{!}{%
\begin{tabular}{>{\centering\bfseries}p{0.6cm}
                >{\raggedright}p{7.2cm}
                >{\raggedright\arraybackslash}p{5.2cm}}
\toprule
\textbf{ID} &
\textbf{Objective} &
\textbf{Target\_Metric} \\
\midrule
O1 & 13-agent CrewAI architecture; 5-state FSM covering 56 LFS fields
     (Sections A--K), 3 employment paths, 11 skip gates &
     All skip gates correct; $<$3 s P95 per turn \\
O2 & 4-stage hierarchical ISCO-08 RAG (441 unit groups, 4 Qdrant
     collections, weighted confidence, Claude re-ranking) &
     $\geq$85\% top-1, $\geq$92\% top-3, $\kappa \geq 0.80$ \\
O3 & Full 4-level ISIC Rev.~4 industry classification
     (Section $\to$ Division $\to$ Group $\to$ 4-digit Class) &
     $\geq$80\% top-1 at Class level \\
O4 & Dual ISCED: ISCED~2011 attainment (0--8) + ISCED-F~2013
     field (Broad $\to$ Narrow $\to$ 4-digit Detailed) &
     $\geq$85\% level; $\geq$75\% 4-digit field \\
O5 & \textbf{Semantic Relation Engine}: ISCO$\leftrightarrow$ISIC$\leftrightarrow$ISCED
     three-way crosswalk using ILO/UNESCO tables;
     SemanticCoherence score 0--1; HITL on HIGH violations &
     Score $\geq$0.90 on coherent cases; all HIGH violations escalated \\
O6 & 5-language support (en/ar/ar-dialect/ur/hi/tl); 30-rule Arabic dialect
     normalisation; code-switch detection; bilingual NER &
     Lang-ID $\geq$98\%; CSW F1 $\geq$92\% \\
O7 & Person Register pre-fill (40--50\% question reduction);
     ILO ICLS-19 rules R01--R10; email OTP + JWT auth &
     40--50\% Q reduction; all 10 rules enforced \\
O8 & GDPR audit trail; bilingual EN+AR report; HITL supervisor
     dashboard; Docker 5-service deployment &
     Zero audit gaps; 10-yr retention; $\leq$17 min median session \\
\bottomrule
\end{tabular}
}
\end{table}

% ------------------------------------------------------------
\section{Research Contributions}
\label{sec:contributions}
% ------------------------------------------------------------

\noindent\textbf{C1 --- Semantic Relation Engine (primary novel contribution).}
A deterministic three-way crosswalk between ISCO-08, ISIC Rev.~4, and
ISCED~2011 that produces a calibrated SemanticCoherence score (0--1)
using the ILO Geneva 2012 ISCO-ISIC table and UNESCO ISCED~2011 Operational
Manual Table~7.
Severity bands: None ($\geq$0.90, $+10\%$ confidence adjustment), LOW
($\geq$0.70, $+5\%$), MODERATE ($\geq$0.50, $-5\%$), HIGH ($<$0.50, $-20\%$,
HITL mandatory).
Sub-major-group stricter rules apply (e.g.\ sub-major 22 Health Professionals:
ISCED $\geq$7 required; sub-major 25 ICT Professionals: ISIC section J
preferred).

\noindent\textbf{C2 --- Complete LFS in 5 languages.}
Full 56-field Sections A--K questionnaire (ILO ICLS-19) with 3 dynamic
employment paths, 11 conditional skip gates, and 2 ILO-mandated re-routing
rules implemented in English, Arabic MSA+regional Arabic dialect, Urdu, Hindi, and
Tagalog --- the first complete multilingual implementation of the LFS.

\noindent\textbf{C3 --- Four-stage hierarchical ISCO-08 RAG (441 unit groups).}
Parent-code payload-filtered pipeline across four Qdrant collections using
weighted confidence ($c = 0.10s_1 + 0.20s_2 + 0.20s_3 + 0.50s_4$),
ISCO-08 synonym enrichment (50+ unit-group descriptions), keyword pre-filter
for major-group anchoring, and LLM re-ranking skipped when top similarity
$\geq$0.92.

\noindent\textbf{C4 --- Simultaneous multi-standard classification.}
The first conversational survey system to classify each respondent's occupation
(ISCO-08, 4-digit), industry (ISIC Rev.~4, 4-digit class), education attainment
(ISCED~2011, 0--8), and education field (ISCED-F~2013, 4-digit detailed) within
a single session, with cross-standard coherence validation.

\noindent\textbf{C5 --- Production-grade open-source deployment.}
A Dockerised 5-service stack (backend/FastAPI, frontend/Next.js~14, PostgreSQL~15,
Qdrant, Redis~7) with Email OTP + JWT HS256 authentication, 11-table PostgreSQL
schema, GDPR Art.~15 compliant immutable audit logging, and bilingual EN+AR
employment report generation.

% ------------------------------------------------------------
\section{Scope and Limitations}
\label{sec:scope}
% ------------------------------------------------------------

The empirical validation is scoped to the LFS in multilingual GCC contexts.
Tagalog support reflects the study region's Filipino community; Bengali (from the
original proposal) was replaced by Tagalog based on demographic data from the study region.

The person-register integration uses data from previous simulated survey rounds.
Live integration with the national identity authority
register is outside project scope but the modular architecture accommodates it.

The ISIC Rev.~4 classifier covers 100+ 4-digit class entries; full coverage of
all 238 classes with equal depth requires a dedicated corpus not currently
available.

The pilot study uses college-student volunteers aged 18--25; results may be
more favourable than in the general LFS population.

% ------------------------------------------------------------
\section{Thesis Organisation}
\label{sec:organisation}
% ------------------------------------------------------------

\noindent\textbf{Chapter~\ref{chap:relatedwork}} surveys six research areas, identifies
eight critical gaps, and positions the Semantic Relation Engine within prior
literature on automated occupation coding, RAG, multilingual NLP, and AI in
official statistics.

\noindent\textbf{Chapter~\ref{chap:architecture}} presents the 13-agent CrewAI
architecture, the 5-state FSM, the 4-stage hierarchical ISCO-08 RAG pipeline,
the ISIC Rev.~4 and ISCED classifiers, and the Semantic Relation Engine in full
technical detail.

\noindent\textbf{Chapter~\ref{chap:implementation}} describes the Docker deployment,
the 11-table PostgreSQL schema, Qdrant vector collections, Arabic dialect
normalisation, ILO ICLS-19 validation rules (R01--R10), and the bilingual
Next.js frontend with RTL Arabic support.

\noindent\textbf{Chapter~\ref{chap:evaluation}} specifies the three-system benchmark
(BM25 / Flat RAG / Hierarchical RAG), the 100-case evaluation set, the Semantic
Relation Engine 10-case demo, and the randomised pilot study protocol.

\noindent\textbf{Chapter~\ref{chap:results}} reports all results and confirms that all
twelve performance targets were met.

\noindent\textbf{Chapter~\ref{chap:conclusion}} synthesises contributions, limitations,
and seven future research directions.
